\documentclass[11pt]{article}
\usepackage[a4paper,margin=1.9cm]{geometry}
\usepackage[T1]{fontenc}
\usepackage{textcomp}
\usepackage[spanish]{babel}
\usepackage{amsmath,amssymb}
\usepackage{lmodern}
\renewcommand{\familydefault}{\sfdefault}
\usepackage{enumitem}

\newcommand{\vect}[2]{\begin{pmatrix}#1\\#2\end{pmatrix}}
\newcommand{\vectiii}[3]{\begin{pmatrix}#1\\#2\\#3\end{pmatrix}}
\usepackage{tikz}
\usetikzlibrary{calc,arrows.meta,patterns,decorations.pathmorphing}
\usepackage[table]{xcolor}
\usepackage{multicol}
\usepackage{parskip}

\definecolor{unalblue}{RGB}{31,73,124}

\newlist{opciones}{enumerate}{1}
\setlist[opciones]{label=(\alph*),itemsep=6pt,topsep=6pt,leftmargin=1.8em,font=\normalfont}
\setlist[enumerate,1]{itemsep=1.4em,topsep=1em}

\newcommand{\examheader}{%
  \noindent
  \begin{minipage}[t]{0.6\linewidth}
    {\small\bfseries Geometría Vectorial y Analítica --- Parcial 3}\\
    {\small Semestre 2016--01}
  \end{minipage}%
  \hfill
  \begin{minipage}[t]{0.35\linewidth}
    \raggedleft\small Escuela de Matemáticas. Universidad Nacional de Colombia, Sede Medellín.
  \end{minipage}\\[2pt]
  \rule{\linewidth}{0.6pt}
}
\newcommand{\examfooter}{%
  \vfill
  \rule{\linewidth}{0.4pt}\\[1pt]
  {\scriptsize\textit{Transcripción digital --- MathUNAL}\hfill\textcolor{unalblue}{\textbf{\thepage}}}
}
\pagestyle{empty}

\newcommand{\nota}[1]{{\small\itshape\color{unalblue!70!black} #1}}

\begin{document}
\examheader

\begin{center}
{\large Tercer Examen Parcial de Geometría}\\[4pt]
{\small 12 de marzo del 2016}
\end{center}

\vspace{6pt}
\textbf{Puntaje. Sólo para uso Oficial}

\begin{center}
\begin{tabular}{|c|c|c|c|c|c|}
\hline
1 (/30) & 2 (/30) & 3 (/10) & 4-6 (/30) & Total & Nota \\
\hline
 & & & & & \\
\hline
\end{tabular}
\end{center}

\vspace{10pt}
\textbf{Instrucciones:} Antes de empezar verifique que su examen esté completo y que sus datos sean correctos. Por favor verifique que su celular esté apagado. No está permitido el uso de calculadora ni de otros aparatos electrónicos. Solucione las preguntas en los espacios en blanco asignados a cada una de ellas. No está permitido sacar hojas en blanco ni ningún tipo de apuntes durante el examen. \textbf{Duración: 1 hora y 50 minutos.}

\vspace{6pt}
\nota{En cada uno de los literales de los puntos 1 y 2, debe presentar detalladamente el procedimiento para llegar a la solución. NOTA: Respuestas sin procedimiento no se tendrán en cuenta.}

\begin{enumerate}
\item Considere la recta $\mathcal{L}$ dada por la ecuación general $3x+y+10=0$.

\begin{enumerate}[label=(\alph*)]
\item{}[10pt] Encuentre unas ecuaciones escalares paramétricas para la recta $\mathcal{L}$.
\item{}[10pt] Considere la recta $\mathcal{L}'$ que es perpendicular a $\mathcal{L}$ y pasa por el punto $Q=\vect{2}{2}$. Encuentre el punto de intersección entre $\mathcal{L}$ y $\mathcal{L}'$.
\item{}[10pt] Encuentre el punto $P$ de la recta $\mathcal{L}$, que está más cercano al origen.
\end{enumerate}

\item Considere la recta $\mathcal{L}$ dada por la ecuación general $\sqrt{3}x+y+\sqrt{3}=0$.

\begin{enumerate}[label=(\alph*)]
\item{}[10pt] Encuentre el ángulo de inclinación de $\mathcal{L}$.
\item{}[10pt] Encuentre la ecuación de la bisectriz del ángulo obtuso formado por $\mathcal{L}$ y el eje $x$. \textbf{Observación}: Recuerde que la bisectriz de un ángulo es la recta que divide al ángulo en dos ángulos iguales.
\item{}[10pt] Halle dos vectores unitarios, normales a la recta $\mathcal{L}$.
\end{enumerate}

\item{}[10pt] Sean $V$ y $W$ dos vectores no nulos y ortogonales. Demuestre que $\|V+W\|=\|V-W\|$.
\end{enumerate}

\vspace{6pt}
\nota{En las preguntas 4 a 6 complete los espacios en blanco o elija la (las) opción (opciones) correcta (correctas) según el caso. Tenga presente que los datos en cada literal son independientes y pueden variar. NOTA: En esta sección se califica sólo la respuesta y no se tiene en cuenta el procedimiento.}

\begin{enumerate}
\setcounter{enumi}{3}
\item{}[10pt] En la figura se muestra la recta $\mathcal{L}$, con pendiente $m$, que pasa por el origen, así como también los vectores $U$, $V$ y $W$. Responda a las siguientes preguntas.

\begin{center}
\begin{tikzpicture}[scale=0.9]
\draw[-{Latex}] (-3,0) -- (3,0) node[right]{$x$};
\draw[-{Latex}] (0,-2.5) -- (0,2.5) node[above]{$y$};
\draw[thick] (-2,2) -- (2,-2) node[below right]{$\mathcal{L}$};
\draw[-{Latex},thick] (0,0) -- (1.2,3) node[above]{$V=\vect{v_1}{3}$};
\draw[-{Latex},thick] (0,0) -- (1,-2) node[right]{$U=\vect{1}{-2}$};
\draw[-{Latex},thick] (0,0) -- (-1,2.4) node[above left]{$W=\vect{-1}{w_2}$};
\end{tikzpicture}
\end{center}

\begin{enumerate}[label=(\alph*)]
\item{}[3pt] Si $\|W\|=\sqrt{10}$ entonces $w_2=$ \underline{\hspace{3cm}}.
\item{}[3pt] Si $U$ es normal a $\mathcal{L}$ entonces $m=$ \underline{\hspace{3cm}}.
\item{}[4pt] Si $U\cdot V=-2$ entonces $\|V\|=$ \underline{\hspace{3cm}}.
\end{enumerate}

\item{}[10pt] Las rectas $\mathcal{L}_1$, $\mathcal{L}_2$, $\mathcal{L}_3$ de la figura tienen ángulos de inclinación $\alpha_1,\alpha_2,\alpha_3$ y pendientes $m_1,m_2,m_3$ respectivamente. Se entiende que $\mathcal{L}_1$ y $\mathcal{L}_2$ son perpendiculares.

\begin{center}
\begin{tikzpicture}[scale=0.9]
\draw[-{Latex}] (-3,0) -- (3.5,0) node[right]{$x$};
\draw[-{Latex}] (0,-1.5) -- (0,2.5) node[above]{$y$};
\draw[thick] (-2.2,-1.6) -- (2,1.5) node[right]{$\mathcal{L}_2$};
\draw[thick] (-1,-1.6) -- (1.6,1.9) node[above]{};
\draw[thick] (-3,1.2) -- (3,-0.9);
\node at (2.6,-0.75) {$\mathcal{L}_1$};
\node at (0.9,-0.15) {$\mathcal{L}_3$};
\node[font=\small] at (-1.4,-0.15) {$\alpha_1$};
\node[font=\small] at (0.3,0.35) {$\theta$};
\node[font=\small] at (0.55,-0.15) {$\alpha_3$};
\node[font=\small] at (1.5,-0.15) {$\alpha_2$};
\end{tikzpicture}
\end{center}

\begin{enumerate}[label=(\alph*)]
\item{}[3pt] Si $m_1=\dfrac{1}{2}$ entonces $\tan\alpha_2=$ \underline{\hspace{3cm}}.
\item{}[4pt] Si $\mathcal{L}_2: 3x+2y-9=0$ y además $\mathcal{L}_3=\left\{t\vect{1}{3}+\vect{2}{0}: t\in\mathbb{R}\right\}$, entonces $\cos\theta=$ \underline{\hspace{3cm}}.
\item{}[3pt] Si $\tan\alpha_1=1$ y $m_3=3$ entonces $\tan\theta=$ \underline{\hspace{3cm}}.
\end{enumerate}

\nota{Indicación. Se recuerda que $\tan(\beta+\gamma)=\dfrac{\tan\beta+\tan\gamma}{1-\tan\beta\tan\gamma}$ y $\tan(\beta-\gamma)=\dfrac{\tan\beta-\tan\gamma}{1+\tan\beta\tan\gamma}$.}

\item{}[10pt] En la figura a continuación, las rectas $\mathcal{L}_1$ y $\mathcal{L}_3$ son ortogonales, el vector $N$ es su punto de intersección y dirige a $\mathcal{L}_3$. La recta $\mathcal{L}_2$ pasa por el origen y es paralela a $\mathcal{L}_1$, $U$ es un vector del plano.

\begin{center}
\begin{tikzpicture}[scale=0.9]
\draw[-{Latex}] (-3,0) -- (3,0) node[right]{$x$};
\draw[-{Latex}] (0,-1.5) -- (0,2.5) node[left]{$y$};
\draw[thick] (-2.5,-1.2) -- (2,2.4) node[above]{$\mathcal{L}_3$};
\draw[thick] (-2.7,-1.9) -- (2.5,1.5) node[right]{$\mathcal{L}_1$};
\draw[thick] (-2.7,1.35) -- (2.7,-1.35) node[below right]{$\mathcal{L}_2$};
\draw[-{Latex},thick] (0,0) -- (0.9,1.9) node[right]{$U=\vect{u_1}{u_2}$};
\node[left,font=\small] at (-0.3,0.15) {$N=\vect{-2}{2}$};
\end{tikzpicture}
\end{center}

\begin{enumerate}[label=(\alph*)]
\item{}[3pt] Si $\mathrm{P}_{\mathcal{L}_3}U=\vect{-4}{4}$ entonces $\operatorname{dist}(U,\mathcal{L}_1)=$ \underline{\hspace{3cm}}.
\item{}[3pt] Si $\|\mathrm{P}_{\mathcal{L}_2}U\|=2$ y $\|U\|=4$ entonces $\operatorname{dist}(U,\mathcal{L}_2)=$ \underline{\hspace{3cm}}.
\item{}[4pt] Suponga que $Y=t\vect{2}{-2}+U$ con $t\in\mathbb{R}$ y además $\operatorname{dist}(Y,\mathcal{L}_3)=3$ entonces $\operatorname{dist}(U,\mathcal{L}_3)=$ \underline{\hspace{3cm}}.
\end{enumerate}

\end{enumerate}

\examfooter
\end{document}
